 \documentclass[12pt, a4paper, american]{article}
\usepackage[a4paper]{geometry}
%\newgeometry{left=1in, bottom=1in, top=1in, right=1in}
%\newgeometry{left=1.25in, bottom=1.25in, top=1.25in, right=1.25in}
\renewcommand{\baselinestretch}{1.5}
\usepackage{amsmath,a4,babel,booktabs,graphicx,rotating,float,amsfonts,mathtools, amsthm, amssymb}
\usepackage[LGR,T1]{fontenc} % spell out all text encodings used
\usepackage[utf8]{inputenc} %
\usepackage{caption}
\usepackage{subcaption} % for subfigures 
\usepackage{graphicx, mathptmx}
%\usepackage[position=top]{subfig}
\usepackage[sort&compress]{natbib}

\usepackage[plain]{fancyref} %fancy references
\usepackage{microtype,framed} %other packages added
\usepackage{wasysym}
\usepackage{multirow,array}
\usepackage{multicol}
\usepackage{tcolorbox} % to create colored text boxes
\usepackage{bbm}
\usepackage{graphicx}
\usepackage{enumitem} 
\usepackage{xcolor}
\usepackage{abbrevs}
\usepackage{pgffor}
\usepackage{comment}
\usepackage{nicefrac}
\setlength{\parskip}{5pt}
\usepackage{setspace}
\usepackage{titling}
\onehalfspacing 
\usepackage{pdflscape}
\usepackage{threeparttable}
\usepackage{soul}
\usepackage{clipboard}
\newclipboard{myclipboard}
\usepackage[citecolor= black,
			colorlinks=true,
		    linkcolor = black,
            urlcolor  = black]{hyperref}
\usepackage{xurl}
\usepackage{chngcntr}   % for per-chapter numbering of figures
\usepackage{float}



\newenvironment{q}%for the appendix
{
\setlength{\parindent}{0pt}
}
%
\makeatletter
\usepackage{tikz}
\usepackage{verbatim}
\usetikzlibrary{positioning, shapes, arrows, fit, backgrounds}
\newcommand*{\h}{\hspace{5pt}}% for indentation
\usepackage{epigraph} % for nice quote 
\usepackage{makecell} % for nice table cells

\newcommand{\ssection}[1]{\section[#1]{\centering #1}}

\newgeometry{top=2.2cm, bottom=2.2cm, left=2.2cm, right=2.2cm}

% ------------ Mathematical environments -----------

\theoremstyle{plain}
\newtheorem{theorem}{Theorem}

\theoremstyle{definition}
\newtheorem{prediction}[theorem]{Prediction}


% *****************************************************************
% Estout related things
% *****************************************************************
\newcommand{\sym}[1]{\rlap{#1}}% Thanks to David Carlisle

\let\estinput=\input% define a new input command so that we can still flatten the document

\newcommand{\estwide}[3]{
		\vspace{.75ex}{
			\begin{tabular*}
			{\textwidth}{@{\hskip\tabcolsep\extracolsep\fill}l*{#2}{#3}}
			\toprule
			\estinput{#1}
			\bottomrule
			\addlinespace[2.50ex]
			\end{tabular*}
			}
		}	

\newcommand{\estauto}[3]{
		\vspace{.75ex}{
			\begin{tabular}{l*{#2}{#3}}
			\toprule
			\estinput{#1}
			\bottomrule
			\addlinespace[2.75ex]
			\end{tabular}
			}
		}

% Allow line breaks with \\ in specialcells
	\newcommand{\specialcell}[2][c]{%
	\begin{tabular}[#1]{@{}c@{}}#2\end{tabular}}


%% For quote

\makeatletter
\newenvironment{chapquote}[2][3em]
  {\setlength{\@tempdima}{#1}%
   \def\chapquote@author{#2}%
   \parshape 1 \@tempdima \dimexpr\textwidth-2\@tempdima\relax%
   \itshape}
  {\par\normalfont\hfill--\ \chapquote@author\hspace*{\@tempdima}\par\bigskip}
\makeatother

% *****************************************************************


% *****************************************************************
% Custom subcaptions
% *****************************************************************
% Note/Source/Text after Tables

% Note/Source/Text after Tables
\newcommand{\tabtext}[1]{%
	\vspace{-0.5ex}
 \begin{tablenotes}[para,flushleft]
 \hspace{5pt}
 \hangindent=1.75em
 #1
 \end{tablenotes}
}

\newcommand{\tabnote}[1]{\parbox{\linewidth}{\tabtext{\emph{Note:~}~#1}}}
\newcommand{\tabsource}[1]{\tabtext{\emph{Source:~}~#1}}
\newcommand{\starnote}{\vspace{0.75ex}\tabtext{* $p$\;<\;0.10, ** $p$\;<\;0.05, *** $p$\;<\;0.01. Robust standard errors in parentheses.}}
\newcommand{\starnotebootstrap}{\vspace{0.75ex}\tabtext{* $p$\;<\;0.10, ** $p$\;<\;0.05, *** $p$\;<\;0.01. Bootstrapped standard errors in parentheses.}}

\newcommand{\starnotesimple}{\vspace{0.75ex}\tabtext{* $p$\;<\;0.10, ** $p$\;<\;0.05, *** $p$\;<\;0.01.}}

\newcommand{\starnotenonrobust}{\tabtext{* $p$\;<\;0.10, ** $p$\;<\;0.05, *** $p$\;<\;0.01. Standard errors in parentheses.}}

% *****************************************************************
% Custom subcaptions: Figure
% *****************************************************************
% Note/Source/Text after Tables
\newcommand{\figtext}[1]{
	\vspace{-0.5ex}
	\captionsetup{justification=justified,font=small}
	\caption*{\hspace{8pt}\hangindent=1.5em #1}
	}
%\newcommand{\fignote}[1]{\figtext{\scriptsize Notes:~#1}}
\newcommand{\fignote}[1]{\figtext{\emph{Note:~}~#1}}

\newcommand{\figsource}[1]{\figtext{\emph{Source:~}~#1}}

% Character substitution that prints brackets and the minus symbol in text mode and does not reserve any space. Thanks to David Carlisle
\def\yyy{%
  \bgroup\uccode`\~\expandafter`\string-%
  \uppercase{\egroup\edef~{\noexpand\text{\llap{\textendash}\relax}}}%
  \mathcode\expandafter`\string-"8000 }

\def\xxxl#1{%
\bgroup\uccode`\~\expandafter`\string#1%
\uppercase{\egroup\edef~{\noexpand\text{\noexpand\llap{\string#1}}}}%
\mathcode\expandafter`\string#1"8000 }

\def\xxxr#1{%
\bgroup\uccode`\~\expandafter`\string#1%
\uppercase{\egroup\edef~{\noexpand\text{\noexpand\rlap{\string#1}}}}%
\mathcode\expandafter`\string#1"8000 }
\def\textsymbols{\xxxl[\xxxr]\xxxl(\xxxr)\yyy}

% *****************************************************************
% siunitx
% *****************************************************************
\usepackage{siunitx} % centering in tables
	\sisetup{
		detect-mode,
		tight-spacing		= true,
		group-digits		= false ,
		input-signs		= ,
		input-symbols		= ( ) [ ] - + *,
		input-open-uncertainty	= ,
		input-close-uncertainty	= ,
		table-align-text-post	= false
        }

\usepackage{footmisc} %NR: Changed this to use numeric footnotes again
\DefineFNsymbolsTM{myfnsymbols}{% def. from footmisc.sty "bringhurst" symbols
  \textasteriskcentered *
  \textdagger    \dagger
  \textdaggerdbl \ddagger
  \textsection   \mathsectionwe 
%  \textbardbl    \|%
%  \textparagraph \mathparagraph
}%
\setfnsymbol{myfnsymbols}



\theoremstyle{definition}
\newtheorem{prop}[theorem]{Proposition}
\newtheorem{corollary}[theorem]{Corollary}
%\newtheorem{prediction}[theorem]{Prediction}
\newtheorem{hypothesis}[theorem]{Hypothesis}

\newtheorem{result}{Result}



% ----------------------- Abbreviations ----------------------------------

% General abbreviations
\newcommand{\lucid}{Lucid}
\newcommand{\fox}{\textit{Fox News}}
\newcommand{\msnbc}{\textit{MSNBC}}
\newcommand{\bidenrescueplan}{\textit{Biden Rescue Plan}}
\newcommand{\americanrescueplan}{\textit{American Rescue Plan}}

\newcommand{\onlineappendix}{Online Appendix}

% Statistical abbreviations
\newcommand{\n}{$n$\;=\;}
\newcommand\onestar{$p$\;<\;0.10}
\newcommand\twostar{$p$\;<\;0.05}
\newcommand\threestar{$p$\;<\;0.01}
\newcommand\nostar{$p$\;>\;0.10}
\newcommand\pt{$p$\;<\;0.001}
\newcommand{\p}{$p$\;=\;}

% Experiment Instructions
\newcommand{\includeInstructions}[1]{
	\begin{center}
		\includegraphics[width=0.90\textwidth]{#1.png}
	\end{center}
}


\newcommand{\includeInstructionsLong}[1]{
	\begin{center}
		\includegraphics[height=0.95\textheight]{#1.png}
	\end{center}
}






